% Source: https://artofproblemsolving.com/community/c3h2302477_ab2024s_math_hard_questions
% Stable identity: aops-topic-2302477-q10
\textbf{Modulo $3$.} Since $4\equiv1\pmod3$ and $n^3\equiv n\pmod 3$ (Fermat's little theorem),
\[4n^4-n\equiv n^4-n=n\left(n^3-1\right)\equiv n(n-1)\pmod 3,\]
so $3\mid 4n^4-n$ exactly when $n\equiv0$ or $1\pmod3$.

\textbf{Modulo $5$.} If $5\mid n$ then $n^4-1\equiv-1\pmod 5$; otherwise $n^4\equiv1\pmod5$ by Fermat's little theorem. So $5\mid n^4-1$ exactly when $5\nmid n$.

\textbf{Counting.} Whether both conditions hold depends only on $n$ modulo $15$. Of the $15$ residues, those with $n\equiv0,1\pmod 3$ and $n\not\equiv0\pmod5$ number $2\cdot4=8$ (by the Chinese remainder theorem each pair of residues modulo $3$ and $5$ occurs exactly once). Since $15\mid1089!$, the set $\{1,\dots,1089!\}$ consists of $\dfrac{1089!}{15}$ complete blocks of $15$ consecutive integers, so the count is
\[\boxed{\dfrac{8\cdot1089!}{15}}.\]
